% Options for packages loaded elsewhere
% Options for packages loaded elsewhere
\PassOptionsToPackage{unicode}{hyperref}
\PassOptionsToPackage{hyphens}{url}
\PassOptionsToPackage{dvipsnames,svgnames,x11names}{xcolor}
%
\documentclass[
  letterpaper,
  DIV=11,
  numbers=noendperiod]{scrartcl}
\usepackage{xcolor}
\usepackage{amsmath,amssymb}
\setcounter{secnumdepth}{-\maxdimen} % remove section numbering
\usepackage{iftex}
\ifPDFTeX
  \usepackage[T1]{fontenc}
  \usepackage[utf8]{inputenc}
  \usepackage{textcomp} % provide euro and other symbols
\else % if luatex or xetex
  \usepackage{unicode-math} % this also loads fontspec
  \defaultfontfeatures{Scale=MatchLowercase}
  \defaultfontfeatures[\rmfamily]{Ligatures=TeX,Scale=1}
\fi
\usepackage{lmodern}
\ifPDFTeX\else
  % xetex/luatex font selection
\fi
% Use upquote if available, for straight quotes in verbatim environments
\IfFileExists{upquote.sty}{\usepackage{upquote}}{}
\IfFileExists{microtype.sty}{% use microtype if available
  \usepackage[]{microtype}
  \UseMicrotypeSet[protrusion]{basicmath} % disable protrusion for tt fonts
}{}
\makeatletter
\@ifundefined{KOMAClassName}{% if non-KOMA class
  \IfFileExists{parskip.sty}{%
    \usepackage{parskip}
  }{% else
    \setlength{\parindent}{0pt}
    \setlength{\parskip}{6pt plus 2pt minus 1pt}}
}{% if KOMA class
  \KOMAoptions{parskip=half}}
\makeatother
% Make \paragraph and \subparagraph free-standing
\makeatletter
\ifx\paragraph\undefined\else
  \let\oldparagraph\paragraph
  \renewcommand{\paragraph}{
    \@ifstar
      \xxxParagraphStar
      \xxxParagraphNoStar
  }
  \newcommand{\xxxParagraphStar}[1]{\oldparagraph*{#1}\mbox{}}
  \newcommand{\xxxParagraphNoStar}[1]{\oldparagraph{#1}\mbox{}}
\fi
\ifx\subparagraph\undefined\else
  \let\oldsubparagraph\subparagraph
  \renewcommand{\subparagraph}{
    \@ifstar
      \xxxSubParagraphStar
      \xxxSubParagraphNoStar
  }
  \newcommand{\xxxSubParagraphStar}[1]{\oldsubparagraph*{#1}\mbox{}}
  \newcommand{\xxxSubParagraphNoStar}[1]{\oldsubparagraph{#1}\mbox{}}
\fi
\makeatother


\usepackage{longtable,booktabs,array}
\usepackage{calc} % for calculating minipage widths
% Correct order of tables after \paragraph or \subparagraph
\usepackage{etoolbox}
\makeatletter
\patchcmd\longtable{\par}{\if@noskipsec\mbox{}\fi\par}{}{}
\makeatother
% Allow footnotes in longtable head/foot
\IfFileExists{footnotehyper.sty}{\usepackage{footnotehyper}}{\usepackage{footnote}}
\makesavenoteenv{longtable}
\usepackage{graphicx}
\makeatletter
\newsavebox\pandoc@box
\newcommand*\pandocbounded[1]{% scales image to fit in text height/width
  \sbox\pandoc@box{#1}%
  \Gscale@div\@tempa{\textheight}{\dimexpr\ht\pandoc@box+\dp\pandoc@box\relax}%
  \Gscale@div\@tempb{\linewidth}{\wd\pandoc@box}%
  \ifdim\@tempb\p@<\@tempa\p@\let\@tempa\@tempb\fi% select the smaller of both
  \ifdim\@tempa\p@<\p@\scalebox{\@tempa}{\usebox\pandoc@box}%
  \else\usebox{\pandoc@box}%
  \fi%
}
% Set default figure placement to htbp
\def\fps@figure{htbp}
\makeatother


% definitions for citeproc citations
\NewDocumentCommand\citeproctext{}{}
\NewDocumentCommand\citeproc{mm}{%
  \begingroup\def\citeproctext{#2}\cite{#1}\endgroup}
\makeatletter
 % allow citations to break across lines
 \let\@cite@ofmt\@firstofone
 % avoid brackets around text for \cite:
 \def\@biblabel#1{}
 \def\@cite#1#2{{#1\if@tempswa , #2\fi}}
\makeatother
\newlength{\cslhangindent}
\setlength{\cslhangindent}{1.5em}
\newlength{\csllabelwidth}
\setlength{\csllabelwidth}{3em}
\newenvironment{CSLReferences}[2] % #1 hanging-indent, #2 entry-spacing
 {\begin{list}{}{%
  \setlength{\itemindent}{0pt}
  \setlength{\leftmargin}{0pt}
  \setlength{\parsep}{0pt}
  % turn on hanging indent if param 1 is 1
  \ifodd #1
   \setlength{\leftmargin}{\cslhangindent}
   \setlength{\itemindent}{-1\cslhangindent}
  \fi
  % set entry spacing
  \setlength{\itemsep}{#2\baselineskip}}}
 {\end{list}}
\usepackage{calc}
\newcommand{\CSLBlock}[1]{\hfill\break\parbox[t]{\linewidth}{\strut\ignorespaces#1\strut}}
\newcommand{\CSLLeftMargin}[1]{\parbox[t]{\csllabelwidth}{\strut#1\strut}}
\newcommand{\CSLRightInline}[1]{\parbox[t]{\linewidth - \csllabelwidth}{\strut#1\strut}}
\newcommand{\CSLIndent}[1]{\hspace{\cslhangindent}#1}



\setlength{\emergencystretch}{3em} % prevent overfull lines

\providecommand{\tightlist}{%
  \setlength{\itemsep}{0pt}\setlength{\parskip}{0pt}}






\KOMAoption{captions}{tableheading}
\makeatletter
\@ifpackageloaded{caption}{}{\usepackage{caption}}
\AtBeginDocument{%
\ifdefined\contentsname
  \renewcommand*\contentsname{Table of contents}
\else
  \newcommand\contentsname{Table of contents}
\fi
\ifdefined\listfigurename
  \renewcommand*\listfigurename{List of Figures}
\else
  \newcommand\listfigurename{List of Figures}
\fi
\ifdefined\listtablename
  \renewcommand*\listtablename{List of Tables}
\else
  \newcommand\listtablename{List of Tables}
\fi
\ifdefined\figurename
  \renewcommand*\figurename{Figure}
\else
  \newcommand\figurename{Figure}
\fi
\ifdefined\tablename
  \renewcommand*\tablename{Table}
\else
  \newcommand\tablename{Table}
\fi
}
\@ifpackageloaded{float}{}{\usepackage{float}}
\floatstyle{ruled}
\@ifundefined{c@chapter}{\newfloat{codelisting}{h}{lop}}{\newfloat{codelisting}{h}{lop}[chapter]}
\floatname{codelisting}{Listing}
\newcommand*\listoflistings{\listof{codelisting}{List of Listings}}
\makeatother
\makeatletter
\makeatother
\makeatletter
\@ifpackageloaded{caption}{}{\usepackage{caption}}
\@ifpackageloaded{subcaption}{}{\usepackage{subcaption}}
\makeatother
\usepackage{bookmark}
\IfFileExists{xurl.sty}{\usepackage{xurl}}{} % add URL line breaks if available
\urlstyle{same}
\hypersetup{
  pdftitle={Separating Anthropogenic Forcing from Solar-Locked Natural Variability in Global Ocean SST Anomalies (1850--2026)},
  pdfauthor={Dr.~Alf Loeffler},
  pdfkeywords={sea surface temperature, singular spectrum
analysis, solar cycles, greenhouse gas forcing, coupled
oscillators, Richards growth curve},
  colorlinks=true,
  linkcolor={blue},
  filecolor={Maroon},
  citecolor={Blue},
  urlcolor={Blue},
  pdfcreator={LaTeX via pandoc}}


\title{Separating Anthropogenic Forcing from Solar-Locked Natural
Variability in Global Ocean SST Anomalies (1850--2026)}
\usepackage{etoolbox}
\makeatletter
\providecommand{\subtitle}[1]{% add subtitle to \maketitle
  \apptocmd{\@title}{\par {\large #1 \par}}{}{}
}
\makeatother
\subtitle{A Non-Parametric Subspace, Coupled-Oscillator, and Sigmoidal
Growth Framework}
\author{Dr.~Alf Loeffler}
\date{2026-11-05}
\begin{document}
\maketitle


\textbf{Keywords:} Sea surface temperature; singular spectrum analysis;
solar cycles; greenhouse gas forcing; coupled oscillators; Richards
growth curve

\subsection{Abstract}\label{abstract}

Global ocean heat content accounts for the vast majority of Earth's
surface thermal capacity. Anthropogenic greenhouse gas emissions disturb
the equilibrium of incoming solar and outgoing longwave radiation; while
both components are attenuated, the retention of outgoing thermal
radiation dominates, driving persistent heat content growth.

Mathematically, ocean thermal dynamics can be modeled as a system of
coupled oscillators driven by multi-decadal solar irradiance cycles.
Observed sea surface temperature (SST) anomalies integrate these
long-term dynamics alongside short-term seasonal and meteorological
``noise'', necessitating a clear separation of secular trends from
internal periodic variations.

Using NOAA ERSSTv6 global area-weighted sea surface temperature (SST)
anomalies (1850--2026), we extract background secular trends using three
complementary approaches: (1) solar-cycle phase locking using splines,
(2) non-parametric Singular Spectrum Analysis (SSA), and (3) non-linear
logistic growth modeling. Rolling-correlation analysis reveals distinct
structural breaks matching the early acceleration phase of the modern
times rise of the SST. These breaks mark a mid-century transition regime
where anthropogenic warming begins to obscure natural periodic
variability.

Our findings demonstrate that historical multi-decadal
hiatuses---specifically around \textasciitilde1910 and between 1940 and
1970---are driven by destructive phase alignments of natural cycles
rather than shifts in external forcing Medhaug et al. (2017). While
sigmoidal growth curves successfully isolate the monotonic anthropogenic
trajectory, post-inflection data through 2026 remain insufficient to
establish a hard physical upper asymptote for ocean heat uptake.

To evaluate the human contribution, long-term natural periodic
variability is separated from the underlying anthropogenic trend.

\subsection{Introduction}\label{introduction}

The global ocean acts as a massive thermal reservoir absorbing the solar
forced heating. Thermodynamic laws are the principal forces trying to
obtain equilibrium between the various ocean basins. Huge atmospheric
and oceanic currents transport the heat and are observed as
oscillations. Some of these are stable for centuries others change with
seasonal or daily periods.

Mathematically, large-scale oceanic thermal adjustments can be
conceptualized through coupled-oscillator frameworks. Analogous to the
`recharge oscillator' mechanics developed for equatorial ocean dynamics
by Burgers, Jin, and Oldenborgh (2005) , global sea surface temperature
anomalies integrate delayed subsurface and boundary feedback responses,
which manifest as prominent multi-decadal periods when stripped of
secular trends.

This physical system can be conceptualized through three main dynamics:

\begin{enumerate}
\def\labelenumi{\arabic{enumi}.}
\tightlist
\item
  \emph{Coupled Solar Driving:} The ocean behaves as a set of non-linear
  coupled oscillators driven by periodic variations in solar irradiance
  across Schwabe (\textasciitilde11-year), Hale
  (\textasciitilde22-year), and Gleissberg (\textasciitilde88-year)
  cycles.
\item
  \emph{Radiative Imbalance:} Anthropogenic emissions of long-lived
  greenhouse gases alter the planetary energy balance primarily by
  retarding outgoing longwave thermal radiation emitted primary in the
  northern hemisphere but distribute with a time lag of
  \textasciitilde{} 2 years across the full global surface
  (\(4\pi R^2\)). While local atmospheric processes also affect incoming
  solar irradiance, the net global energetic disturbance is
  overwhelmingly dominated by the retention of outgoing thermal
  radiation over the total ocean area.
\item
  \emph{Signal Disentanglement:} Isolating the anthropogenic warming
  trend requires separating the underlying secular trajectory from
  natural, multi-decadal oscillatory modes.
\end{enumerate}

Traditional Fourier-based spectral filtering and linear detrending needs
to be combined with a separation of the dominant periods
(\emph{principal component analysis}) to understand the effect of the
periodical varying solar input.

In this study, we apply splines to lock the anomaly to solar irradiance,
paired with non-parametric Singular Spectrum Analysis (SSA) and
sigmoidal growth curve modeling to isolate natural cyclical variations
from anthropogenic warming without enforcing artificial linearity or
spectral symmetry.

\subsection{Data and Pre-Processing}\label{data-and-pre-processing}

\subsubsection{NOAA ERSSTv6 Global Ocean
Anomalies}\label{noaa-ersstv6-global-ocean-anomalies}

Global monthly sea surface temperature (SST) fields are obtained from
the NOAA Extended Reconstructed Sea Surface Temperature version 6
(ERSSTv6) covering January 1850 through 2026 (Huang et al. 2025).
Spatial operations on the \(180 \times 90\) global grid
(\(2^\circ \times 2^\circ\) resolution) were executed using the
\texttt{terra\ package} (Hijmans et al. 2026). Global ocean mean
anomalies were generated using an area-weighted ocean mask based on
Berkeley Earth grid boundaries (Rohde and Hausfather 2020). Annual and
semi-annual seasonal signals were removed using the
\texttt{itsmr\ package} (Weigt 2025).

To account for high-frequency vortex and weather noise, both mean and
median anomaly metrics were extracted across the time series. As shown
with Figure~\ref{fig-Global-SST-Anomaly} the median anomaly series
exhibits a lower overall slope and a narrower dynamic range. Beyond
2000, the lower bound of the median series approaches that of the mean
anomaly, reflecting the physical lower limit of seawater freezing points
versus nearly unconstrained upper atmospheric warming spikes. For
primary modeling, we prioritize the area-weighted mean SST anomaly, as
it strictly conserves enthalpy and total thermal energy, while using the
median to verify against localized polar amplification artifacts.
Although we attempt to analyze the change of the global ocean heat
content, we take SST as a proxy because of the available data processed
and published by NOAA Huang et al. (2025)

\begin{figure}

\centering{

\includegraphics[width=0.85\linewidth,height=\textheight,keepaspectratio]{Figs_Ocean/Fig1_Mean_Median.png}

}

\caption{\label{fig-Global-SST-Anomaly}Global Temperature Anomaly (Mean
vs Median)}

\end{figure}%

\subsubsection{Solar Irradiance Proxy
(1749--2026)}\label{solar-irradiance-proxy-17492026}

A very extensive analysis of the solar irradiation is published by (Lean
2018) covering paleontologic times until today. We use a simplification
of the method as input data for our analysis covering the time from 1850
to 2026. The model evaluates long-term solar activity using daily
sunspot numbers from WDC SILSO v2.0 (Clette and Lefèvre 2015). From the
smoothed time series irradiance variations are calculated using
satellite-calibrated quadratic polynomial transformations (Ambelu et al.
2011).

\begin{figure}

\centering{

\includegraphics[width=0.85\linewidth,height=\textheight,keepaspectratio]{Figs_Ocean/Fig2_Solar_trend.png}

}

\caption{\label{fig-sol-power}Solar Power Periodicity}

\end{figure}%

The solar power series is decomposed using SSA to map long-term
modulation envelopes, including the Gleissberg - (\textasciitilde88
years), Hale cycle (\textasciitilde20 years) and Schwabe cycle.The
Schwabe cycle's \textasciitilde11-year period is transformed into
\textasciitilde22 years as power is approximated by a 2nd-degree
polynomial. The irradiation variance calculated from the Schwabe cycle
shown as the blue line and Hale cycle as red line periods.

\subsection{Physical Framework \& Signal
Decomposition}\label{physical-framework-signal-decomposition}

The Figure~\ref{fig-Global-SST-Anomaly} of the global ocean anomaly
suggests a strong upwards trend around \textasciitilde{} 1950. A
standard methods to calculate this decomposition is an approximation
with polynomials or repeated differentiation Brockwell and Davis (2016).
Both methods gave questionable results with the ocean anomaly time
series: Polynomials need to made a decision on the degree used for
fitting and the result depends on this. Differentiation produces high
frequencies with the data complicating further analysis. In addition
extended polynomials violate the physical boundaries of ocean anomalies.

Locking the SST anomaly to the periodic changing solar irradiation gives
physical correct model. The second method \emph{Singular Spectrum
Analysis} Golyandina, Nekrutkin, and Zhigljavsky (2001) includes the
extraction of the dominant harmonic periods by principal component
analysis (PCA).

Both methods include no physical derived criterion or method to separate
the effect of greenhouse gases from the natural periodic changes of the
anomaly.This is achieved by regression to a mathematical model that
respects the physical constraints of SST. A lower limit ( the freezing
point of seawater) and an upper limit caused by increased evaporation
and as consequence cloud formation. Such a limit of\(\approx 35\,^\circ\text{C}\)we
observe in some partial ocean areas.

\subsubsection{Solar-Locked Baseline and Natural Periodic
Changes}\label{solar-locked-baseline-and-natural-periodic-changes}

By connecting the observed SST to the solar cycle minima, a
\textbf{solar-locked baseline} is established. Subtracting this baseline
yields a stationary residual vector containing multi-decadal periodic
variability without monotonic secular trends.

A LOESS smoothing fit (\(\text{span} = 0.25\), \(\approx 44\) years)
(Cleveland and Devlin 1988) isolates the background trajectory
(sec.base;green line). The difference between the solar-locked trend and
the secular baseline highlights a primary \textasciitilde22-year
oscillation corresponding to the Schwabe power cycle and the
\textasciitilde20-year Hale cycle.

\begin{figure}

\centering{

\includegraphics[width=0.85\linewidth,height=\textheight,keepaspectratio]{Figs_Ocean/Fig3_sec_sol.png}

}

\caption{\label{fig-sol-lckd-trend}Solar-Locked / Secular Trend}

\end{figure}%

\subsubsection{Singular Spectrum Analysis
(SSA)}\label{singular-spectrum-analysis-ssa}

The dominant periods, are extracted from the observed SST anomalies
using Rssa package (Golyandina and Korobeynikov 2014). An embedding
window of \(L = 703\) months (\(\approx 58.6\) years, or
\(\approx N/3\)) was selected to cleanly resolve multi-decadal modes
without distorting boundary values.

\begin{figure}

\centering{

\includegraphics[width=0.85\linewidth,height=\textheight,keepaspectratio]{Figs_Ocean/Fig4_SSA_sol.png}

}

\caption{\label{fig-ssa-decomp}SSA Decomposition of Solar-Locked
Components}

\end{figure}%

The ssa.trend consists of the first two components, the harmonics
represent the components 3 to 20 . Comparing Figure~\ref{fig-ssa-decomp}
with the Figure~\ref{fig-sol-lckd-trend} shows that the main harmonics
calculated with SSA look very similar to the difference of the solar
locked trend and the secular reference. The intuitive solar locked
method gives very similar results to the mathematical derived
decomposition. The SSA method can be applied also to the residuals of
the anomaly providing a consistent method for ocean analysis.

\subsection{Anthropogenic Trend \& Growth
Modeling}\label{anthropogenic-trend-growth-modeling}

\subsubsection{Logistic and Richards Growth Curve
Fitting}\label{logistic-and-richards-growth-curve-fitting}

Assuming limits on SST the secular anthropogenic component is modeled
using sigmoidal growth functions. The initial effect of the irradiation
unbalance on the SST is observed as an upward trend expanding the
historical stable range of the anomaly. A mathematical model must also
respect an upper limit asymptotical. An upper limit is clearly observed
with the ``West Pacific Warm Pool'' a partial basins of the global
ocean. A sigmoidal function can map both. Hence as a very simplified
model we fit a standard symmetric logistic curve:

\begin{equation}\phantomsection\label{eq-logistic}{
y(t) = b + \frac{L}{1 + e^{-k(t - t_0)}}
}\end{equation}

An asymmetric Generalized Logistic (Richards) growth curve could
conceptually capture the upper and lower limit different. However,
fitting the Richards curve failed due to matrix singularities during
optimization. This numerical failure points to a deeper physical cause:
the global anomaly time series aggregates vast spatial areas, thereby
diluting the sharp regional feedback mechanisms---such as the regional
recharge/discharge limits observed in smaller sub-basins---that are
required to mathematically constrain an additional asymmetry parameter.

\begin{equation}\phantomsection\label{eq-richards}{
y(t) = L_r + \frac{\Delta}{\left(1 + e^{-k(t - t_{0r})}\right)^{1/\nu}}
}\end{equation}

where \(t_0\) represents the inflection point, \(k\) is the growth rate,
and \(L\) denotes the upper asymptotic limit.

\begin{figure}

\centering{

\includegraphics[width=0.85\linewidth,height=\textheight,keepaspectratio]{Figs_Ocean/Fig5_Industr_separated.png}

}

\caption{\label{fig-sigmoid-ssa-low}Sigmoidal Growth Fitting to Low-Pass
SSA Reconstruction}

\end{figure}%

Figure~\ref{fig-sigmoid-ssa-low} shows the lowpass filtered series
(observed,green line); the ``industrial trend'' (red line); the
inflection point of the sigmoid: dot at \(t_0 \approx 2018\); and the
residuals (pre-ind, blue-line)

With data extending to 2026, post-inflection observations (\(\sim 8\)
years) provide initial curvature constraints but are insufficient to
establish the exact physical asymptote (\(L\)). In this context, the
fitted limit reflects short-term mixed-layer heat capacity and thermal
inertia rather than a final ceiling on global climate change.

\subsubsection{Structural Break \& Natural Residual
Analysis}\label{structural-break-natural-residual-analysis}

Subtracting the sigmoid growth model Figure~\ref{fig-sigmoid-ssa-low}
(industr) from the total anomaly (observed) yields the
\textbf{pre-industrial natural residual} (pre-ind). An 11-year rolling
correlation between the observed anomaly and natural reconstruction
shows high correlation until brk1 \(\sim 1950\), where a clear
structural break indicates the transition into an
anthropogenic-dominated climate regime.

\begin{figure}

\centering{

\includegraphics[width=0.85\linewidth,height=\textheight,keepaspectratio]{Figs_Ocean/Fig6_roll_cor.png}

}

\caption{\label{fig-rolling-cor}Rolling Correlation showing the
Mid-Century Structural Break}

\end{figure}%

The two vertical lines Figure~\ref{fig-rolling-cor} (brk1 and brk2
breaks of the correlation) mark the transition period and are also shown
in Figure~\ref{fig-sigmoid-ssa-low}.

\subsection{Discussion \& Results}\label{discussion-results}

\subsubsection{Logistic Inflection 2009 Phase Anchor, and Asymptotic
Constraints}\label{logistic-inflection-2009-phase-anchor-and-asymptotic-constraints}

A critical question in interpreting the non-linear trajectory is the
timing of the logistic curve's inflection point. While unconstrained
logistic regression places the symmetry point near 2018 , the natural
harmonic minimum is firmly anchored near 2009 Figure~\ref{fig-nat-har}.
By restricting the logistic regression through the 2009 phase minimum
and verifying that recalculated natural residuals remain virtually
identical, we confirm an iterative, physically consistent correction
robust against circular logic. Correcting the turning point to 2009
naturally lowers the model's fitted upper asymptote (\(L\)).
Nevertheless, this adjusted asymptote reflects short-term mixed-layer
heat capacity rather than a definitive physical ceiling on global
warming. A full accounting of upper thermal bounds requires integrating
regional regime-shift dynamics and boundary-layer feedbacks highlighted
by structural break analysis.

\begin{figure}

\centering{

\includegraphics[width=0.85\linewidth,height=\textheight,keepaspectratio]{Figs_Ocean/Fig7_nat_decomposed.png}

}

\caption{\label{fig-natrl-ssa-decmp}Natural Harmonics Decomposed}

\end{figure}%

\subsubsection{Key Findings}\label{key-findings}

By strictly separating solar-locked natural cycles from the secular
warming trend, this analysis provides a physical baseline for ocean heat
uptake that avoids polynomial overfitting and empirical distortion.

\begin{figure}

\centering{

\includegraphics[width=0.85\linewidth,height=\textheight,keepaspectratio]{Figs_Ocean/Fig8_nat_summary.png}

}

\caption{\label{fig-nat-har}Natural Harmonics Discussion}

\end{figure}%

\subsubsection{Structural Breaks and the 1940--1970
Hiatus}\label{structural-breaks-and-the-19401970-hiatus}

During the 1940--1970 hiatus, accelerating greenhouse warming was
temporarily offset by the downward phase of natural multi-decadal
harmonics, masking underlying secular acceleration.

\begin{itemize}
\tightlist
\item
  \textbf{Hiatus Dynamics:} The mid-20th-century warming slowdown
  (1950--1970) is primarily explained by destructive phase interference
  between natural solar harmonics (notably the \textasciitilde22-year
  Hale cycle) and ocean thermal lag, rather than purely external aerosol
  forcing.
\item
  \textbf{Secular Trajectory:} Isolating the anthropogenic mode via SSA
  demonstrates a continuous, monotonic warming trajectory that aligns
  with cumulative planetary energy imbalance and thermodynamic enthalpy
  conservation.
\item
  \textbf{Regime Stability:} Structural break testing confirms that
  multi-decadal oscillations act as superposed periodic variations
  around a single, evolving non-linear warming baseline.
\end{itemize}

\subsubsection{Future Work}\label{future-work}

This diagnostic framework will be extended to basin-specific regional
SST anomalies (Atlantic Multidecadal Oscillation and Pacific Decadal
Oscillation dynamics) to evaluate local resonance effects and regional
thermal inertia differences.

\subsubsection{Statements and
Declarations}\label{statements-and-declarations}

Funding: The author declares that no funds, grants, or other support
were received during the preparation of this manuscript.

Competing Interests: The author has no competing interests to declare
that are relevant to the content of this article.

Data Availability: Global monthly sea surface temperature anomaly fields
are publicly available from the NOAA Extended Reconstructed Sea Surface
Temperature version 6 (ERSSTv6) repository at https://www.ncei.noaa.gov.
Daily sunspot numbers are available from WDC SILSO v2.0 at
https://sidc.be/silso/.

Code Availability: The complete analysis pipeline, Quarto source
documents, and R scripts supporting the findings of this study are
openly archived in the public GitHub repository
https://github.com/aloes2512/Ocean\_Analysis.

Corresponding author: alf.loeffler@icloud.com

\subsubsection{Acknowledgments}\label{acknowledgments}

Artificial Intelligence assistance (Google Gemini) was utilized strictly
as an auxiliary tool for verifying R processing scripts, optimizing
Quarto syntax formatting, and refining text grammar. The mathematical
conceptualization, thermal dynamics framework, data analysis, and
scientific conclusions were developed independently by the author.

\phantomsection\label{refs}
\begin{CSLReferences}{1}{0}
\bibitem[\citeproctext]{ref-ambelu2011estimation}
Ambelu, T., E. O. Falayi, E. O. Elemo, and O. Oladosu. 2011.
{``Estimation of Total Solar Irradiance from Sunspot Number.''}
\emph{Latin-American Journal of Physics Education} 5 (2): 374--81.

\bibitem[\citeproctext]{ref-brockwell2016introduction}
Brockwell, Peter J., and Richard A. Davis. 2016. \emph{Introduction to
Time Series and Forecasting}. 3rd ed. New York: Springer.
\url{https://doi.org/10.1007/978-3-319-29854-2}.

\bibitem[\citeproctext]{ref-burgers2005simplest}
Burgers, Geert, Fei-Fei Jin, and Geert Jan van Oldenborgh. 2005. {``The
Simplest ENSO Recharge Oscillator.''} \emph{Geophysical Research
Letters} 32 (13): L13706. \url{https://doi.org/10.1029/2005GL022951}.

\bibitem[\citeproctext]{ref-SILSOSunspotCounts2015}
Clette, F., and L. Lefèvre. 2015. {``{SILSO Sunspot Number V2.0}.''}
https://doi.org/10.24414/qnza-ac80.
\url{https://doi.org/10.24414/qnza-ac80}.

\bibitem[\citeproctext]{ref-cleveland1988}
Cleveland, William S., and Susan J. Devlin. 1988. {``Locally Weighted
Regression: An Approach to Regression Analysis by Local Fitting.''}
\emph{Journal of the American Statistical Association} 83 (403):
596--610. \url{https://doi.org/10.1080/01621459.1988.10478639}.

\bibitem[\citeproctext]{ref-golyandina2014rssa}
Golyandina, Nina, and Anton Korobeynikov. 2014. {``Basic Singular
Spectrum Analysis and Forecasting with r: {Rssa}.''} \emph{Journal of
Statistical Software} 60 (9): 1--25.
\url{https://doi.org/10.18637/jss.v060.i09}.

\bibitem[\citeproctext]{ref-golyandina2001}
Golyandina, Nina, Vladimir Nekrutkin, and Anatoly Zhigljavsky. 2001.
\emph{Analysis of Time Series Structure: SSA and Related Techniques}.
Boca Raton, FL: Chapman; Hall/CRC.

\bibitem[\citeproctext]{ref-Hijmans2026}
Hijmans, Robert J., Andrew Brown, A. Márcia Barbosa, Emanuele Cordano,
and Krzysztof Dyba. 2026. \emph{Terra: Spatial Data Analysis}.
\url{https://github.com/rspatial/terra}.

\bibitem[\citeproctext]{ref-huang2025ersstv6}
Huang, Boyin, Chunying Liu, Viva Banzon, Eric Freeman, Richard Graham,
Bill Hankins, Thomas M. Smith, and Huai-Min Zhang. 2025. {``Extended
Reconstructed Sea Surface Temperature, Version 6 (ERSSTv6): Upgraded
Reconstruction and QC Methods.''} \emph{Journal of Climate} 38 (4):
1125--43. \url{https://doi.org/10.1175/JCLI-D-24-0123.1}.

\bibitem[\citeproctext]{ref-lean2018estimating}
Lean, J. L. 2018. {``Estimating Solar Irradiance Since 850 {CE}.''}
\emph{Earth and Space Science} 5 (4): 133--49.
\url{https://doi.org/10.1002/2017EA000357}.

\bibitem[\citeproctext]{ref-medhaug2017reconciling}
Medhaug, Iselin, Martin B. Stolpe, Erich M. Fischer, and Reto Knutti.
2017. {``Reconciling Controversies about the 'Global Warming Hiatus'.''}
\emph{Nature} 545 (7652): 41--47.
\url{https://doi.org/10.1038/nature22315}.

\bibitem[\citeproctext]{ref-rohde2020berkeley}
Rohde, Robert A., and Zeke Hausfather. 2020. {``The Berkeley Earth
Land/Ocean Temperature Record.''} \emph{Earth System Science Data} 12
(10): 3469--79. \url{https://doi.org/10.5194/essd-12-3469-2020}.

\bibitem[\citeproctext]{ref-Weigt2025}
Weigt, George. 2025. \emph{Itsmr: Time Series Analysis Using the
Innovations Algorithm}.
\url{https://doi.org/10.32614/CRAN.package.itsmr}.

\end{CSLReferences}




\end{document}
